% Part: lambda-calculus
% Chapter: introduction
% Section: term-revisited

\documentclass[../../../include/open-logic-section]{subfiles}

\begin{document}

\olfileid{lam}{syn}{tr}
\olsection{పదాలను $\alpha$-తుల్యతా వర్గాలుగా చూడడం}

ఇకపై పదాలను $\alpha$-తుల్యత వరకు పరిగణిస్తాం. అంటే ఒక
పదాన్ని రాసినప్పుడు, అది ఉన్న $\alpha$-తుల్యతా వర్గాన్ని
ఉద్దేశిస్తాం. ఉదాహరణకు, $\lambd[a][\lambd[b][a c]]$ అని
రాస్తే, $\lambd[a][\lambd[b][a c]]$,
$\lambd[b][\lambd[a][b c]]$ మొదలైన, దానికి
$\alpha$-తుల్యమైన అన్ని పదాల సమితిని ఉద్దేశిస్తాం.

మునుపటి విభాగాల్లో $N, Q$ వంటి అక్షరాలు ఒక్కొక్క పదాన్ని
సూచించేవి. ఇకపై అవి ఒక్కొక్క తుల్యతా వర్గాన్ని సూచిస్తాయి;
తరువాతి అధ్యయనంలో పదాల బదులు ఈ వర్గాలే మన అంశాలు.
$x, y$ వంటి అక్షరాలు మాత్రం చరాలనే సూచిస్తాయి.

వర్గం~$M$లోని ఏదైనా \tetoken{మూలకాన్ని}{!!{element}} సూచించడానికి
$\rep{M}$ అనే సంకేతాన్ని వాడతాం. ఒకటి కంటే ఎక్కువ కావాలంటే
$\rep{M}[0], \rep{M}[1], etc. $ అనే మూల సంకేతాలను వాడతాం.

వివరణ సులభంగా ఉండటానికి పదాలపై వాడిన సంకేతాలనే మళ్లీ వాడతాం.
వర్గాలపై కింది నిర్వచనాలు ఇస్తాం:
\begin{defn}
  \begin{enumerate}
  \item $\lambd[x][N]$ను $\lambd[x][\rep{N}]$ను కలిగి ఉన్న
    వర్గంగా నిర్వచిస్తాం.
  \item $PQ$ను $\rep{P}\rep{Q}$ను కలిగి ఉన్న వర్గంగా నిర్వచిస్తాం.
  \end{enumerate}
\end{defn}

ఇవి ప్రతినిధిని ఎంచుకున్న విధానంపై ఆధారపడవు:
$\alpha$-పరివర్తనం అమూర్తీకరణ, ప్రయోగాలతో అనుకూలంగా ఉంటుంది.

\begin{defn} \ollabel{def:fv}
  $\alpha$-తుల్యతా వర్గం~$M$ యొక్క \emph{స్వేచ్ఛా చరాలు},
  లేదా $FV(M)$, అనేవి $FV(\rep{M})$గా నిర్వచించబడతాయి.
\end{defn}

\olref[alp]{thm:fv}లో చూపినట్లుగా $FV(\rep{M}[0]) =
FV(\rep{M}[1])$ కనుక ఈ నిర్వచనం ప్రతినిధి ఎంపికపై ఆధారపడదు.

వర్గాల్లో ప్రతిస్థాపనకూ అదే సంకేతాన్ని వాడతాం:
\begin{defn} \ollabel{def:sub}
  $M$లో $y$ స్థానంలో $R$ను \emph{ప్రతిస్థాపించడం}, లేదా
  $\Subst{M}{R}{y}$, అంటే ప్రతిస్థాపన నిర్వచితమయ్యే ఏ
  $\rep{M}$, $\rep{R}$లను తీసుకున్నా
  $\Subst{\rep{M}}{\rep{R}}{y}$ను కలిగి ఉన్న వర్గం.
  \sourcecorrection{OLTELAMTR-001}{మూలం వర్గాలపై $\Subst{M}{R}{y}$ను నిర్వచిస్తానంటూ కుడివైపు పదం $\Subst{\rep{M}}{\rep{R}}{y}$నే రాస్తుంది; అది స్వయంగా వర్గం కాదు. గణిత సంకేతం మార్చకుండా, ఫలితం ఆ పదం ఉన్న $\alpha$-తుల్యతా వర్గమని తెలుగు గద్యంలో స్పష్టంచేశాం.}
\end{defn}

ఈ నిర్వచనం కూడా ప్రతినిధి ఎంపికపై ఆధారపడదని మూలం
\olref[alp]{cor:sub}ను చూపుతుంది.
\sourcecorrection{OLTELAMTR-002}{ఆ ఉపసిద్ధాంతానికి ఆధారమైన $\alpha$-సురక్షిత ప్రతిస్థాపన సిద్ధాంతపు మూల నిరూపణలో OLTELAMALP-005–006 ఖాళీలు ఉన్నాయి. ఇక్కడ ప్రతినిధి-స్వాతంత్ర్యాన్ని కొత్తగా నిరూపించలేదు; మూలం చేసిన ఆధార సూచనను మాత్రమే నిలిపి, పూర్తి నిరూపణ పెండింగ్‌లో ఉందని ప్రకటిస్తున్నాం.}

ఈ నిర్వచనం తర్కాన్ని గణనీయంగా సులభం చేస్తుంది. ఉదాహరణకు:
\begin{align}
  \Subst{\lambd[x][x]}{y}{x} & =\ollabel{eq:1}\\
  &= \Subst{\lambd[z][z]}{y}{x} \ollabel{eq:2}\\
  &= \lambd[z][\Subst{z}{y}{x}] \\
  &= \lambd[z][z]
\end{align}

పదాలపైనే ప్రతిస్థాపనగా చూస్తే \olref{eq:1}
నిర్వచితం కాదు. కానీ ఇప్పుడు మనం వర్గాలపై ప్రతిస్థాపనను
పరిగణిస్తున్నాం. అందుకే \olref{eq:2} సాధ్యం:
$\lambd[x][x]$, $\lambd[z][z]$ ఒకే వర్గానికి చెందినవి
కనుక మొదటిదాని స్థానంలో రెండవదాన్ని ఎంచుకోవచ్చు.

అదే కారణంతో, ఇకపై ప్రతిస్థాపనకు అవసరమైన షరతులు తీరే
ప్రతినిధులనే ఎంచుకున్నట్లు భావిస్తాం. ఉదాహరణకు,
$\Subst{\lambd[x][N]}{R}{y}$ కనిపిస్తే,
$x \neq y$, $x \notin FV(R)$ అయ్యేలా ప్రతినిధి
$\lambd[x][N]$ను ఎంచుకున్నామని భావిస్తాం.

$\lambd[x][x]$ను “వర్గం” అని పిలవడం కొంచెం వింతగా
ఉంటుంది. కాబట్టి పరిచయమైన $\lambda$-పదాల నుంచి వీటిని
వేరుచేయడానికి ఇకపై వీటిని $\Lambda$-పదాలు అంటాం;
ఈ భాగంలో మిగతా చోట్ల సంక్షిప్తంగా “పదాలు” అని కూడా అంటాం.

\begin{editorial}
  పాత పదాలను పూర్తిగా వదిలేయలేం: $\Lambda$-పదాల
  నిర్వచనమంతా $\lambda$-పదాలపైనే ఆధారపడింది. అంతేకాక
  $\Lambda$-పదాలపై ప్రమేయాలను నేరుగా నిర్వచించే విధానాన్ని
  ఇంకా ఇవ్వలేదు. అందువల్ల అలాంటి ప్రమేయాలను ముందుగా
  $\lambda$-పదాలపై నిర్వచించి, ప్రతిస్థాపన విషయంలో చేసినట్లే,
  తరువాత $\Lambda$-పదాలపైకి దింపాలి. అలా $\Lambda$-పదాలపై
  ప్రమేయాలను నిర్వచించవచ్చని పాఠకులు సహజంగా అర్థం చేసుకోగలరని భావిస్తున్నాం.
\end{editorial}
\end{document}
